\begin{claim}{C.7f}{Proved}
Retain the fields, curve, and cyclotomic kernel of C.7e. Fix a finite
connected geometric covering $V\to U_{\bar K}$. There is a finite set
$S_V\subset U(\bar K)$, all of whose points are algebraic over $k_0$,
such that, for every $x\in U(K)\setminus S_V$, every element of
$\Pi^{s_x(H)}$ acts trivially on the geometric fibre of $V$ above $x$.
The same finite set works for every open subgroup of $H$.

For a $k_0$-rational cusp $c$, with parameter $u$ over $k_0$, fix a
finite geometric Galois cover with group $G$ and an inertia subgroup
$I\subset G$ above $c$. There is a finite set
$A_G\subset K^\times$, consisting of elements algebraic over $k_0$,
such that, for the tangent $u=a t$ and every $a\notin A_G$, the
image of the full fixed group in $G$ is exactly $I$.
Again this holds for every open subgroup of $H$.

Thus the exceptional set is finite at each fixed finite quotient.
These assertions do not give a uniform finite exceptional set for
all quotients, and do not decide the point $2$ or direction $1$.
\end{claim}
\begin{proof}
We use the following precise consequence of the characteristic-zero
torsion-packet theorem of Baker and Poonen
\cite[Theorem~2, printed p.~110; proof pp.~111--114]{BC_BakerPoonen2001}:
if a smooth proper curve $D$ has genus at least four, the set
\[
 T_D=\{P\in D(\bar K): [P-Q]\text{ is torsion for some }Q\ne P\}
                                                               \tag{C.7f.1}
\]
is finite. Their theorem gives finitely many nontrivial torsion
packets in these genera, and each packet is finite by the
Manin--Mumford theorem. The genus restriction is essential to this
particular input: genus-two curves, and certain genus-three curves,
have infinitely many nontrivial packets.

Let $L$ be the relative algebraic closure of $k_0$ in $\bar K$.
All the geometric covers in use descend to $L$, as explained in
C.7e. Refine the Galois closure of $V$ to a connected geometric
Galois cover $W\to U_{\bar K}$ whose smooth compactification $D$
has genus at least four. Here is an explicit existence argument.
C.7e constructs a refinement whose compactification has genus at
least two for every hyperbolic type. A connected cyclic degree-three
\emph{unramified} cover of that compactification exists: choose a
nonzero order-three line bundle and its trivialized cube. Its genus
is $3(g-1)+1\ge4$. Take a Galois closure of the resulting open
cover over $U$. The compactification still maps to the genus-at-least-four
curve, and hence has genus at least four by Riemann--Hurwitz.

The finite set $T_D$ lies in $D(L)$. Indeed the torsion points of
$\operatorname{Jac}(D)$ are $L$-rational. For a fixed nonzero torsion
point $b$, the fibre over $b$ of
\[
       D\times D\longrightarrow\operatorname{Jac}(D),\qquad
                       (P,Q)\longmapsto[P-Q]
\]
is finite. One can use Baker--Poonen Proposition~4 directly:
off the diagonal the difference map is an embedding when $D$ is
nonhyperelliptic and is finite of degree two in the hyperelliptic
case. Thus its closed points over $b\in\operatorname{Jac}(D)(L)$
are defined over the algebraically closed field $L$.
Set $S_V$ equal to the image of $T_D\cap W$ in $U$.

Take $x\notin S_V$, an open subgroup $H'\subset H$, and a fixed
element $g\in\Pi^{s_x(H')}$. After a finite extension $M/K$, the
cover, all its deck transformations, and the fibre over $x$ are
defined over $M$, and $H_M=G_{M\Q_p(\mu_\infty)}$ lies in $H'$.
The $H_M$-invariant loop represented by $g$ lifts from any chosen
point $P$ of that fibre to a point $Q$ in the same fibre. It is an
$H_M$-invariant path between the corresponding ordinary point
sections. Filling the boundary and using C.7e's rational abelian
restriction lemma gives $[P-Q]$ torsion in
$\operatorname{Jac}(D)(M)$. If $P\ne Q$, then $P\in T_D$ and
$x\in S_V$, a contradiction. Therefore $g$ fixes every point of
the fibre of $W$, and hence of $V$.

For the tangential assertion, we can likewise refine the given
Galois cover and retain the quotient map to $G$. It suffices to
prove the assertion for the refined group, so write $D$ for its
compactification and keep the name $G$. If a fixed element has
image outside $I$, lift its invariant tangential path starting at
a chosen point $P$ above $c$. Its terminal point lies at a distinct
point $Q$ above $c$. This distinction follows from the usual
identification of points over $c$ with inertia cosets; it is not
an assertion that all tangential directions at the same point are
distinct sections. After a finite local extension, filling the
boundary and applying the abelian restriction lemma gives
$[P-Q]$ torsion.

There are only finitely many ordered pairs of distinct points
$P,Q$ above $c$. For each pair whose difference is torsion, fix
$m=m_{PQ}>0$ and $f_{PQ}\in L(D)^\times$ with
\[
                         \operatorname{div}(f_{PQ})=m(P-Q).
\]
Choose parameters $z,w$ at $P,Q$ over $L$. Since the cover is
Galois, the common ramification index above $c$ is an integer $e$.
Write the leading terms as
\[
 \begin{aligned}
 u&=c_Pz^e(1+\cdots),&
 u&=c_Qw^e(1+\cdots),\\
 f_{PQ}&=d_Pz^m(1+\cdots),&
 f_{PQ}&=d_Qw^{-m}(1+\cdots),
 \end{aligned}
 \qquad c_P,c_Q,d_P,d_Q\in L^\times.
                                                               \tag{C.7f.2}
\]
For a tangent $u=a t$, its lifted leading coefficients satisfy
$\alpha^e=a/c_P$ and $\beta^e=a/c_Q$. The normalized formal
units have unique required roots after these leading coefficients
have been chosen, exactly as in C.7e. Applying $f_{PQ}$ to the
invariant path and using the multiplicative local restriction
lemma gives
\[
                  (d_P/d_Q)(\alpha\beta)^m\in\mu_\infty.
\]
Raise this equality to the $e$th power and substitute the two
leading-coefficient equations. It follows that
\[
 a^{2m}\in C_{PQ}\mu_\infty,
 \qquad
 C_{PQ}=(c_Pc_Q)^m(d_Q/d_P)^e\in L^\times.                    \tag{C.7f.3}
\]
For any fixed $C\in L^\times$ and positive integer $r$, the set
\[
                    \{a\in K^\times:a^r\in C\mu_\infty\}
                                                               \tag{C.7f.4}
\]
is either empty or a single coset of the finite group $\mu(K)$.
To see this, the quotient of the $r$th powers of any two solutions
is a root of unity. Their quotient is therefore itself a root of
unity, and belongs to $K$. Conversely multiplying a solution by
an element of $\mu(K)$ preserves the condition. Every solution is
algebraic over $L$, and so lies in $L$ because $L$ is algebraically
closed. Hence it is algebraic over $k_0$.

Let $A_G$ be the union of the finite sets (C.7f.4) arising from
the finitely many eligible pairs $P,Q$. If $a\notin A_G$, an
invariant tangential path cannot change its underlying point over
$c$, so every full fixed element has image in $I$. Conversely the
full cusp inertia is fixed by every open subgroup of $H$, since
the cyclotomic character is trivial there; its image is all of
$I$. This proves the equality.

All exceptional sets were chosen from the fixed geometric cover
and its leading terms. The finite local extension used to contain
the endpoints and the leading roots can vary with $a$, $x$, and
$H'$, but the sets themselves do not. Passing to all finite
quotients gives a countable union of finite sets, not a single
finite exceptional set. The last distinction is exactly why this
result supplies no conclusion at the specified algebraic
parameters.
\end{proof}
